% !TeX root = ../papers/nc-p3.tex
\section{有限代数の三次小行列式と十九次全体の排除}
\label{sec:sectional-nineteen-stein-cubics}
\begin{theorem}\label{thm:sectional-nineteen-stein-cubics}
被約四辺有限標準正則窓の整な非恒等反対八次では、三行小行列式の定める曲線の
重複度込みの総次数は十八以下である。一般点での階数、曲線の被約性、成分数を仮定しない。
\end{theorem}
\begin{proof}[有限代数の次数一の生成元]
十九次と仮定し、定理\ref{thm:sectional-nineteen-model}の基点のブローアップを
$\Gamma$、入力射を $f$、入力束と核束を $H,L$ とする。$R$ を入力四変数の多項式環、
$B=\bigoplus_t H^0(\Gamma,tH)$ を有限直像代数 $f_*\mathcal O_\Gamma$ の全切断加群とする。
一般入力平面切断は前定理の五次曲線 $J$ 自身である。
実際局所完全交叉の基点 $(a,b)$ で一般切断を $a$ と取れば、
ブローアップの $aY=bX$ における対応超平面切断は基点を含めて $a=0$ に同型である。
$H^1(\Gamma,\mathcal O_\Gamma)=0$ である。基点の上のブローアップは局所的に
自明な射影直線束内の相対一次方程式なので、構造層の直像は元の構造層で高次直像は零である。
$J$ の算術種数は二、$\deg H|_J=8$。切断完全列を零から正負両方向に使えば
$H^1(\Gamma,tH)=0$ が全ての整数 $t$ で従う（負方向の最初は定数切断が全射）。
特に一般一次式で割った $B$ は $J$ の全切断加群となり、
$H^1(\mathbb P^3,(f_*\mathcal O_\Gamma)(t))=0$ もLerayの単射から従う。
よって $B$ は次元三のCohen--Macaulay加群である。

一般平面の入力三切断を $W$、$H|_J=P_0$ と書く。
$W\otimes H^0(P_0)\to H^0(2P_0)$ は全射である。評価核二束を $E$ とすると
$\det E=-P_0$、従って余核の双対は $H^0(E\otimes\omega_J)$ である。
この空間は $W\otimes H^0(\omega_J)\to H^0(P_0\omega_J)$ の核。
標準系は基点なしの鉛筆なので、非零な関係があれば $W$ はある
$0\ne s\in H^0(P_0\omega_J^{-1})$ に対し $s H^0(\omega_J)$ を含む
（鉛筆の二生成元の関係はKoszul関係である）。
すると二入力直線が $J$ 上で $s$ の長さ六の零因子を共有する。
しかし $J\to E_H$ の全繊維長は三以下である。実際一般平面上の余核は
$J$ 上の可逆層の有限押出しで、各繊維長は $4-\operatorname{rank}N$ に一致する。
以後の次数では $W$ 内の基点なしの鉛筆の列と $H^1(J,tP_0)=0$ ($t\ge1$) で乗法は全射。
従って $B$ は単位と次数一の四追加切断で生成される。ここで $h^0(P_0)=7$ である。
$B$ のHilbert関数は非負次数で $4t^2+3t+1$（$J$ の切断数を順次加える）。
深さとHilbert分子より極小表示は
\[
 0\longrightarrow R(-2)^3\oplus R(-3)^2\xrightarrow{T}
 R\oplus R(-1)^4\longrightarrow B\longrightarrow0 .
\]
行列式は整八次 $D$ の方程式である。初めの三列の下四行を一次式行列 $A$ とする。
一般入力点で $A$ の階数は三である。さもなくば $D$ 上では単位の繊維値が非零だから
$T$ の第一行は残り四行で張られ、初めの三列は $D$ 上でも階数二以下となる。
その三行小行列式の次数は高々四なので全て恒等零となり $\det T\ne0$ に反する。
\end{proof}

\begin{proof}[二十印での評価と縮約の個数]
$M_D=f_*L$ は全入力上で成立する。前定理の有限な開集合での等号の外は有限集合で、
評価包含の右辺はtorsion-free、左辺は深さ二で有限集合を越えて切断が延長するからである。
$f^{-1}(p)$ の全ての曲線が核空間の点に写れば、核超平面で繊維の像を避けるものがある。
固有性により $p$ の近傍上で $L$ は自明となり、そこで $M_D\simeq f_*\mathcal O_\Gamma$。
従って印 $p$ では $T(p)$ の階数は二以下となる。ただし核次数が正で入力に潰れる曲線の像を除く。

そのような例外入力点は高々二つしかない。定理\ref{thm:cofactor-exclude-nineteen}の
交点評価により、候補は滑らかな二次曲面の二族の母線、あるいは錐の解消の母線と負切断である。
負切断と基点上の例外曲線は核次数零なので除ける。
滑らかな場合、次数二の母線の縮約には二単純基点または一二重基点が必要、
次数三の母線では三単純基点が必要である。三本を潰すには
二単純基点ずつを持つ互いに素な二本と、三単純基点を持つ他方の母線が要る。
しかし三本目は各々に横断的に一度しか会わず、共有できる基点は無限近接込みで各一本につき一点以下。
最初の二本が四中心を使い切っており不可能である。錐の母線も二単純基点または一二重基点を要する。
従って少なくとも十八印で $T(p)$ は階数二以下となる。これらを良い印と呼ぶ。

$A$ の符号付き極大三次小行列式ベクトルを $v$ とする。良い印で $v=0$ なので、
三次評価の同型性から成分の張る空間は二次元以下である。一般階数三ゆえ $v\ne0$。
定数行操作で非零成分だけを最初に並べる。二次元なら $v=g(u_1,u_2,0,0)$ と書け、
$u_i$ は互いに素な同次数の独立な式である。$v^t A=0$ であるため、
各列の初めの二成分は $(u_2,-u_1)$ の倍数である。$A$ が一次式で一般階数三なので
$u_1,u_2$ は一次式、$g$ は二次式である。$g$ の外の印は少なくとも六個である
（定理\ref{thm:finite-cubic-separators}の三次評価と、二次式倍の隣辺の完全跡を使う
\S\ref{sec:reducible-octic-factors}の評価）。
そのうち良い印は共通直線 $u_1=u_2=0$ 上にあるが、直線上は高々三印。
例外を二つ足しても五で足りない。
\end{proof}

\begin{proof}[三次曲面の特異印]
従って一次元の場合しかない。基底を選び直せば $A$ の一行は恒等零で、
残りは三行三列 $C$、$g=\det C\ne0$ は十八個以上の良い印を含む。
零行に対応する $T$ の最後の二成分は二次式 $q_1,q_2$ である。
共通因子は整八次 $\det T$ を割ってしまうから、両式は互いに素である。
この二式の共通零点には印が高々十二個しかない。
実際 $q_1 R_1+q_2 R_1$ は八次元の三次式空間であり、全二十印での三次評価は同型である。
良い印で $C$ の階数が二なら、最初の三列は二次元を張り、$T(p)$ の他の列はそこに入る。
零行に着目すれば $q_1(p)=q_2(p)=0$ である。
従って少なくとも六個の良い印では $C$ の階数は一以下で、$g$ は全てそこで特異である。

整三次曲面の孤立特異点は高々四つである。特異曲線があれば二つの特異点を結ぶ直線は
曲面上にあることから、その曲線は直線で、直線の外にも特異点があれば平面を含んでしまう。
直線上は高々三印。よって整三次は不可能である。
相異なる平面と整二次曲面の積なら、各々が含む印は高々七個と十四個で、
十八印の和集合を得るには共通印は三個以下である。特異印は共通印と錐ならその頂点だけで四以下。
相異なる三平面の積でも延べ印数は二十一以下なので、二平面以上に載る印は高々三個。
重複平面がある場合の台は二平面以下で十四印に達するだけである。
全て六特異印の必要条件に反する。
\end{proof}
\path{computations/verify_sectional_nineteen_stein_cubics.py} は有限代数のHilbert分子、
両母線の基点消費および全ての三次曲面型の印数を検算する。
